% !TeX program = pdflatex
% !TeX encoding = UTF-8
\documentclass[11pt,aspectratio=169,t]{beamer}
% Все презентации курса компилируются только pdflatex.
\usepackage{iftex}
\RequirePDFTeX
\usepackage{cmap}
\usepackage[T2A]{fontenc}
\usepackage[utf8]{inputenc}
\usepackage[english,russian]{babel}
\usepackage{PTSerif}
\usepackage{amsmath,amsfonts,amssymb}
\usepackage{graphicx,microtype,anyfontsize,tikz}
\usepackage{ragged2e}
\usetikzlibrary{arrows.meta,calc}

\usetheme{default}
\usefonttheme{serif}
\DeclareMathOperator{\sen}{sen}
\bibliographystyle{apalike}

% Общие сведения; номер лекции, её название и дата задаются в N.tex.
\newcommand{\coursename}{Введение в маневрирование КА}
\author{Чиняев Владимир Николаевич}
\institute{Московский физико-технический институт}

% Книжное оформление с голубым акцентом.
\definecolor{coursepaper}{HTML}{FBF9F3}
\definecolor{courseink}{HTML}{302B28}
\definecolor{courseaccent}{HTML}{187AA5}
\definecolor{coursemuted}{HTML}{6D6560}
\definecolor{courserule}{HTML}{A9CADA}
\definecolor{coursetint}{HTML}{EDF4F6}
\tikzset{
  course diagram/.style={font=\fontsize{9.5}{11.5}\selectfont,
    line cap=round,line join=round,>=Latex},
  orbit line/.style={draw=courseaccent,line width=1pt},
  aux line/.style={draw=courserule,line width=.5pt},
  vector line/.style={draw=courseink,line width=.65pt,->},
  angle line/.style={draw=courseaccent,line width=.8pt,->},
  orbit point/.style={circle,fill=courseaccent,inner sep=1.7pt}
}
\setbeamercolor{normal text}{fg=courseink,bg=coursepaper}
\setbeamercolor{background canvas}{bg=coursepaper}
\setbeamercolor{structure}{fg=courseaccent}
\setbeamercolor{alerted text}{fg=courseaccent}
\setbeamercolor{frametitle}{fg=courseink,bg=coursepaper}
\setbeamercolor{footline}{fg=coursemuted,bg=coursepaper}
\setbeamercolor{caption name}{fg=courseaccent}

% Основной текст — 11/14 pt; подписи к рисункам — 9.5/11.5 pt.
\let\coursenormalsize\normalsize
\renewcommand{\normalsize}{\coursenormalsize\fontsize{11}{14}\selectfont}
\setbeamerfont{normal text}{size*={11pt}{14pt}}
\setbeamerfont{frametitle}{size*={20pt}{23pt},series=\mdseries}
\setbeamerfont{title}{size*={27pt}{31pt},series=\mdseries}
\setbeamerfont{subtitle}{size*={12.8pt}{15pt},series=\mdseries}
\setbeamerfont{author}{size*={10.8pt}{13pt}}
\setbeamerfont{institute}{size*={9.2pt}{11pt}}
\setbeamerfont{date}{size*={9pt}{11pt}}
\setbeamerfont{footline}{size*={6.8pt}{8pt}}
\setbeamerfont{caption}{size*={9.5pt}{11.5pt}}
\setbeamerfont{itemize/enumerate body}{size*={11pt}{14pt}}
\setbeamerfont{itemize/enumerate subbody}{size*={11pt}{14pt}}
\setbeamerfont{itemize/enumerate subsubbody}{size*={11pt}{14pt}}
\AtBeginDocument{\normalsize}

\setbeamersize{text margin left=6mm,text margin right=6mm}
\setbeamertemplate{navigation symbols}{}
\setbeamertemplate{headline}{}
\setbeamertemplate{caption}[numbered]
\setbeamertemplate{itemize items}[circle]

% Стандартные блоки Beamer позволяют группировать текст и формулы.
\setbeamertemplate{blocks}[default]
\setbeamerfont{block title}{size*={11pt}{14pt},series=\bfseries}
\setbeamerfont{block body}{size*={11pt}{14pt}}
\setbeamercolor{block title}{fg=courseaccent,bg=coursepaper}
\setbeamercolor{block body}{fg=courseink,bg=coursetint}
\setbeamercolor{block title alerted}{parent=block title}
\setbeamercolor{block body alerted}{parent=block body}
\setbeamercolor{block title example}{parent=block title}
\setbeamercolor{block body example}{parent=block body}

% Компактный заголовок; его высота учитывается при размещении содержания.
\setbeamertemplate{frametitle}{%
  \vskip4.5mm%
  \begin{beamercolorbox}[wd=\textwidth,sep=0pt]{frametitle}
    \usebeamerfont{frametitle}\strut\insertframetitle\strut\par
    \vskip2mm%
    {\color{courseaccent}\hrule height .55pt}
  \end{beamercolorbox}%
}

% Нижний колонтитул: курс слева, лекция и номер слайда справа.
\setbeamertemplate{footline}{%
  \begin{beamercolorbox}[wd=\paperwidth,ht=5.5mm,dp=0pt]{footline}
    \begin{tikzpicture}[remember picture,overlay,x=1mm,y=1mm]
      \begin{scope}[shift={(current page.south west)}]
        \draw[courserule,line width=.3pt] (6,5.5)--(154,5.5);
        \node[text=coursemuted,anchor=west,inner sep=0pt,
          font=\usebeamerfont{footline}] at (6,2.8) {\coursename};
        \node[text=coursemuted,anchor=east,inner sep=0pt,
          font=\usebeamerfont{footline}] at (154,2.8)
          {Лекция \arabic{lecture}\quad\textbullet\quad\insertframenumber};
      \end{scope}
    \end{tikzpicture}%
  \end{beamercolorbox}%
}

% Титульник рассчитан на 16:9 (160 × 90 мм), без нижнего колонтитула.
\setbeamertemplate{title page}{%
  \begin{tikzpicture}[remember picture,overlay,x=1mm,y=1mm]
    \begin{scope}[shift={(current page.south west)}]
      \node[text=coursemuted,anchor=north,font=\usebeamerfont{institute}]
        at (80,82) {\insertinstitute};
      \draw[courseaccent,line width=.55pt] (71,71)--(89,71);
      \node[text=courseink,align=center,text width=125mm,
        font=\usebeamerfont{title}] at (80,53)
        {\hyphenpenalty=10000\coursename\par};
      \node[text=courseaccent,align=center,text width=145mm,font=\usebeamerfont{subtitle}]
        at (80,32) {Лекция \arabic{lecture}. \inserttitle};
      \node[text=courseink,font=\usebeamerfont{author}]
        at (80,19) {\insertauthor};
      \node[text=coursemuted,font=\usebeamerfont{date}]
        at (80,9) {\insertdate};
    \end{scope}
  \end{tikzpicture}%
}

\usepackage{booktabs,array}
\setlength{\abovedisplayskip}{5pt}
\setlength{\belowdisplayskip}{5pt}
\setlength{\abovedisplayshortskip}{3pt}
\setlength{\belowdisplayshortskip}{4pt}
\newcommand{\eqcompact}{\fontsize{10}{12}\selectfont}
\newcommand{\figcaption}[1]{{\usebeamerfont{caption}\color{coursemuted}#1\par}}
\newcommand{\dvr}{\Delta\widehat V_r}
\newcommand{\dvt}{\Delta\widehat V_t}
\newcommand{\drh}{\Delta\widehat r}
\newcommand{\vz}{\widehat V_z}
\newcommand{\zh}{\widehat z}
\newcommand{\imp}{\text{\normalfont имп}}
\setcounter{lecture}{5}
\title{Околокруговое приближение.\\Одноимпульсные манёвры. Условия оптимальности.}
\date{20 октября 2026 г.}

\begin{document}

% Только интервалы около формул; основной текст остаётся 11/14 pt.
\newcommand{\derivationmathspacing}{%
  \setlength{\abovedisplayskip}{2pt}%
  \setlength{\belowdisplayskip}{3pt}%
  \setlength{\abovedisplayshortskip}{2pt}%
  \setlength{\belowdisplayshortskip}{2pt}%
  \setlength{\jot}{2pt}}

\begin{frame}[plain]
\titlepage
\end{frame}

% Система граничных условий и материал лекции 4.
\begin{frame}{Полная система граничных условий}
\relax
\derivationmathspacing
\begin{columns}[T,onlytextwidth]
\begin{column}{.64\textwidth}
{\eqcompact\derivationmathspacing
\centering\(\displaystyle
\begin{aligned}[t]
 &2\sum\nolimits_j\Delta\widehat V_{tj}=\Delta a,\\[4pt]
 &\sum\nolimits_j(\Delta\widehat V_{rj}\sin\varphi_j+
              2\Delta\widehat V_{tj}\cos\varphi_j)=\Delta e_x,\\[4pt]
 &\sum\nolimits_j(-\Delta\widehat V_{rj}\cos\varphi_j+
              2\Delta\widehat V_{tj}\sin\varphi_j)=\Delta e_y,\\[4pt]
 &-\sum\nolimits_j\Delta\widehat V_{zj}\sin\varphi_j=0,\\[4pt]
 &\sum\nolimits_j\Delta\widehat V_{zj}\cos\varphi_j=\Delta i,\\[4pt]
 &\sum\nolimits_j\bigl[2\Delta\widehat V_{rj}(1-\cos\varphi_j)
       +\Delta\widehat V_{tj}(-3\varphi_j+4\sin\varphi_j)\bigr]
       =\\[2pt]
 &\hspace{10mm}\Delta t-\frac32\varphi_f\Delta a
       +2\Delta e_x\sin\varphi_f+2\Delta e_y(1-\cos\varphi_f).
\end{aligned}
\)\par}
\end{column}
\begin{column}{.33\textwidth}
{\eqcompact\derivationmathspacing
\[
\begin{aligned}
 \Delta a&=\frac{a_f-a_i}{r_0},\\[2pt]
 \Delta e_x&=e_{xf}-e_{xi},\\[2pt]
 \Delta e_y&=e_{yf}-e_{yi}.
\end{aligned}
\]
}
\RaggedRight
% Баранов, диссертация, с. 38--39: угол между плоскостями орбит.
{\eqcompact\derivationmathspacing
\[
\begin{aligned}
 \cos\Delta\gamma={}&\cos i_f\cos i_i\\
   &+\sin i_f\sin i_i\cos\Delta\Omega,\\[2pt]
 \Delta\Omega={}&\Omega_f-\Omega_i,\\[2pt]
 \Delta i={}&2\sin\frac{\Delta\gamma}{2}\simeq\Delta\gamma.\quad\text{\cite{BaranovDissertation}}
\end{aligned}
\]
}
$i$ --- наклонение; $\Omega$ --- ДВУ.
\par\smallskip
Аргумент широты $\varphi$ отсчитывается от линии узлов.
\end{column}
\end{columns}
\end{frame}

\begin{frame}{Оптимальное импульсное маневрирование}
\relax
При поиске оптимального импульсного перехода необходимо определить \cite{EdelbaumGobetzWashington1966}:
\medskip
\begin{enumerate}
 \setlength{\itemsep}{10pt}
 \item Количество импульсов.
 \item Точки их приложения.
 \item Направления импульсов скорости.
 \item Величины импульсов и распределение суммарной характеристической скорости между ними.
\end{enumerate}
\end{frame}

\begin{frame}{Верхняя оценка числа импульсов}
\relax
{\justifying
В линейной задаче без дополнительных ограничений для оптимального маневрирования достаточно числа импульсов, не превышающего числа независимых граничных условий \cite{Edelbaum1966,Neustadt1964,SternPotter1965}.\par}
\medskip
\begin{center}
\renewcommand{\arraystretch}{1.2}
\begin{tabular}{@{}l@{\qquad}l@{\qquad}c@{}}
\toprule
Задача & Параметры & Число импульсов\\
\midrule
Компланарный переход & $a,e_x,e_y$ & $\leq 3$\\
Пространственный переход & $a,e_x,e_y,i_x,i_y$ & $\leq 5$\\
Компланарная встреча & $a,e_x,e_y,u$ & $\leq 4$\\
Пространственная встреча & $a,e_x,e_y,i_x,i_y,u$ & $\leq 6$\\
\bottomrule
\end{tabular}
\end{center}
\smallskip
Для задачи встречи конечный момент времени задан.
\end{frame}

\begin{frame}{Число импульсов}
\relax
\begin{columns}[T,onlytextwidth]
\begin{column}{.48\textwidth}
\begin{block}{Один импульс}
\begin{itemize}
 \setlength{\itemsep}{5pt}
 \item Переход возможен только между орбитами, имеющими общую точку.
 \item Может уступать нескольким импульсам по затратам характеристической скорости.
 \item Одно включение ДУ упрощает выполнение манёвра и повышает надёжность.
\end{itemize}
\end{block}
\end{column}
\begin{column}{.48\textwidth}
\begin{block}{Несколько импульсов}
\begin{itemize}
 \setlength{\itemsep}{5pt}
 \item Позволяют перейти между орбитами, не имеющими общей точки.
 \item Могут уменьшить суммарную характеристическую скорость.
 \item Повторные включения ДУ предъявляют дополнительные требования к её ресурсу и надёжности.
\end{itemize}
\end{block}
\end{column}
\end{columns}
\end{frame}

\begin{frame}{Изменение большой полуоси}
\relax
\begin{itemize}
 \setlength{\itemsep}{7pt}
 \item Большую полуось изменяет только трансверсальная составляющая импульса.
 \item Величина изменения не зависит от точки приложения импульса.
 \item Радиальная и нормальная составляющие в первом приближении не изменяют большую полуось.
\end{itemize}
\medskip
Величина трансверсальной составляющей:
\[
 2\Delta\widehat V_t=\Delta a,
\]
\[
 \boxed{\Delta\widehat V_t=\frac12\Delta a.}
\]
\end{frame}

\begin{frame}{Изменение вектора эксцентриситета}
\relax
\begin{itemize}
 \setlength{\itemsep}{5pt}
 \item Вектор эксцентриситета изменяют трансверсальная и радиальная составляющие импульса.
 \item Изменение зависит от точки приложения импульса.
\end{itemize}
\[
 \Delta e=\sqrt{(\Delta e_x)^2+(\Delta e_y)^2}.
\]
\begin{columns}[T,onlytextwidth]
\begin{column}{.48\textwidth}
\begin{block}{Трансверсальный импульс}
\centering $\displaystyle |\Delta\widehat V_t|=\frac12\Delta e.$\par\smallskip
\RaggedRight
Одновременно изменяет большую полуось.
\end{block}
\end{column}
\begin{column}{.48\textwidth}
\begin{block}{Радиальный импульс}
\centering $\displaystyle \vphantom{\frac12}|\Delta\widehat V_r|=\Delta e.$\par\smallskip
\RaggedRight
Большая полуось не меняется в первом приближении.
\end{block}
\end{column}
\end{columns}
\end{frame}

\begin{frame}{Изменение эксцентриситета в апсидах}
\relax
{\justifying
Для изменения эксцентриситета при неизменном направлении перицентра трансверсальный импульс прикладывают в перицентре или апоцентре.\par}
\smallskip
\begin{columns}[T,onlytextwidth]
\begin{column}{.48\textwidth}
\begin{block}{В перицентре}
\[
 \Delta\widehat V_t=\frac{e_f-e_i}{2}.
\]
\begin{itemize}
 \setlength{\itemsep}{3pt}
 \item Положительный импульс увеличивает эксцентриситет.
 \item Отрицательный импульс уменьшает эксцентриситет.
\end{itemize}
\end{block}
\end{column}
\begin{column}{.48\textwidth}
\begin{block}{В апоцентре}
\[
 \Delta\widehat V_t=-\frac{e_f-e_i}{2}.
\]
\begin{itemize}
 \setlength{\itemsep}{3pt}
 \item Положительный импульс уменьшает эксцентриситет.
 \item Отрицательный импульс увеличивает эксцентриситет.
\end{itemize}
\end{block}
\end{column}
\end{columns}
\end{frame}

\begin{frame}{Поворот плоскости орбиты}
\relax
\begin{itemize}
 \setlength{\itemsep}{8pt}
 \item Плоскости орбит совмещаются одним нормальным импульсом на линии их пересечения.
 \item Модуль импульса определяется углом $\Delta i$ между плоскостями.
\end{itemize}
\medskip
\[
 \boxed{\Delta\widehat V_z=\pm\Delta i.}
\]
\medskip
\begin{itemize}
 \setlength{\itemsep}{8pt}
 \item Знак «$+$» соответствует восходящему узлу конечной орбиты относительно начальной, знак «$-$» — нисходящему.
 \item В первом приближении большая полуось и эксцентриситет не меняются.
\end{itemize}
\end{frame}

\begin{frame}{Радиальный импульс: орбита}
\relax
\begin{columns}[T,onlytextwidth]
\begin{column}{.50\textwidth}
{\justifying
Исходная орбита — круговая.\par}
\smallskip
\centering% Baranov, Sec. 1.4.1, Fig. 1.9; first-order orbit shape.
% Native size: 70 x 40 mm. Black: reference circle; blue dashed: perturbed orbit.
\begin{tikzpicture}[course diagram,x=1mm,y=1mm]
\path[use as bounding box] (-24,-20) rectangle (46,20);
\draw[vector line] (-18,0)--(30,0) node[above] {$X$};
\draw[vector line] (0,-17)--(0,18) node[left] {$Y$};
\draw[courseink,line width=.65pt] (0,0) circle[radius=13];
\draw[orbit line,dashed] plot[domain=0:360,samples=181,variable=\a]
 ({13*(1+.16*sin(\a))*cos(\a)},{13*(1+.16*sin(\a))*sin(\a)});
\draw[vector line] (0,0)--(130:13) node[midway,left] {$r_0$};
\draw[courseaccent,line width=1.1pt,->] (13,0)--(25,0)
 node[pos=.72,above=1mm] {$\Delta V_r$};
\draw[courseaccent,line width=1.1pt,->] (0,0)--(0,-6)
 node[pos=.7,right=1mm] {$\Delta\boldsymbol e$};
\draw[aux line] (0,15.08)--(16,15.08)--(19,17);
\node[anchor=west,inner sep=0pt] at (19,17) {$\Delta r_\alpha$};
\draw[aux line] (0,-10.92)--(17,-10.92)--(20,-14);
\node[anchor=west,inner sep=0pt] at (20,-14) {$\Delta r_\pi$};
\fill[courseaccent] (13,0) circle[radius=.65];
\node[anchor=north west,inner sep=1mm] at (13,0) {$P$};
\node[anchor=west,text=coursemuted,font=\fontsize{8}{10}\selectfont]
 at (-21,-19) {$\Delta r_\alpha=-\Delta r_\pi=r_0\Delta\widehat V_r>0$};
\end{tikzpicture}

\figcaption{Чёрная линия --- опорный круг; голубой пунктир --- орбита после импульса. Отклонения увеличены.}
\end{column}
\begin{column}{.47\textwidth}
Импульс $\Delta\widehat V_r>0$ приложен при $t=0$, $\theta=n_0t$.
{\eqcompact
\setlength{\abovedisplayskip}{3pt}\setlength{\belowdisplayskip}{3pt}
\setlength{\abovedisplayshortskip}{2pt}\setlength{\belowdisplayshortskip}{2pt}
\[
 \begin{aligned}
 \Delta r&=r_0{\Delta\widehat V_r}\sin\theta,\\
 \Delta a&=0,\\
 \Delta\boldsymbol e&=-{\Delta\widehat V_r}\boldsymbol e_t.
 \end{aligned}
\]
}
Перицентр расположен за $90^\circ$ до точки приложения импульса.
{\eqcompact
\setlength{\abovedisplayskip}{3pt}\setlength{\belowdisplayskip}{3pt}
\setlength{\abovedisplayshortskip}{2pt}\setlength{\belowdisplayshortskip}{2pt}
\[
 \Delta r_\alpha=r_0{\Delta\widehat V_r},\qquad \Delta r_\pi=-r_0{\Delta\widehat V_r}.
\]
}
{\justifying
Здесь $\alpha$ --- апоцентр, $\pi$ --- перицентр.\par}
\end{column}
\end{columns}
\end{frame}

\begin{frame}{Радиальный импульс: фаза}
\relax
\begin{columns}[T,onlytextwidth]
\begin{column}{.50\textwidth}
\centering% Baranov, Sec. 1.4.1, Fig. 1.10. Native size: 70 x 40 mm.
\begin{tikzpicture}[course diagram,x=1mm,y=1mm]
\path[use as bounding box] (-9,-28) rectangle (61,12);
\draw[vector line] (0,-24)--(0,9) node[above right] {$\Delta u$};
\draw[vector line] (0,0)--(57,0) node[above] {$\theta$};
\draw[aux line,dashed] (25,0)--(25,-20);
\draw[orbit line] plot[domain=0:360,samples=121,variable=\a]
 ({\a/7.2},{10*(cos(\a)-1)});
\foreach \x/\lab in {25/\pi,50/2\pi}{
 \draw[courseink,line width=.5pt] (\x,-.8)--(\x,.8);
 \node[above=1mm,inner sep=0pt] at (\x,0) {$\lab$};}
\node[below=1mm,inner sep=0pt] at (25,-20) {$-4\Delta\widehat V_r$};
\fill[courseaccent] (25,-20) circle[radius=.65];
\node[below left,inner sep=1mm] at (0,0) {$0$};
\end{tikzpicture}

\end{column}
\begin{column}{.47\textwidth}
Для $\Delta\widehat V_r>0$:
{\eqcompact
\setlength{\abovedisplayskip}{3pt}\setlength{\belowdisplayskip}{3pt}
\setlength{\abovedisplayshortskip}{2pt}\setlength{\belowdisplayshortskip}{2pt}
\[
 \Delta u=2{\Delta\widehat V_r}(\cos\theta-1).
\]
}
Максимальное отставание при $\theta=\pi$:
{\eqcompact
\setlength{\abovedisplayskip}{3pt}\setlength{\belowdisplayskip}{3pt}
\setlength{\abovedisplayshortskip}{2pt}\setlength{\belowdisplayshortskip}{2pt}
\[
 \Delta u_{\min}=-4\Delta\widehat V_r.
\]
}
Через один опорный период:
{\eqcompact
\setlength{\abovedisplayskip}{3pt}\setlength{\belowdisplayskip}{3pt}
\setlength{\abovedisplayshortskip}{2pt}\setlength{\belowdisplayshortskip}{2pt}
\[
 \Delta u(2\pi)=0,\qquad T_0=\frac{2\pi}{n_0}.
\]
}
{\justifying
Вековое смещение в первом приближении отсутствует.\par}
\end{column}
\end{columns}
\end{frame}

\begin{frame}{Радиальный импульс: траектория}
\relax
Движение относительно КА на опорном круге, $\eta=r_0\Delta u$.
\smallskip
\begin{columns}[T,onlytextwidth]
\begin{column}{.50\textwidth}
\centering% Baranov, Sec. 1.4.1, Fig. 1.11(a). Native size: 70 x 40 mm.
% xi = eta/(r0 q), rho = Delta r/(r0 q), q = Delta V_r/V0 > 0.
\begin{tikzpicture}[course diagram,x=1mm,y=1mm]
\path[use as bounding box] (-59,-20) rectangle (11,20);
\draw[vector line] (-54,0)--(8,0) node[above] {$\eta$};
\draw[vector line] (0,-16)--(0,17) node[above right] {$\Delta r$};
\draw[orbit line] (-24,0) ellipse[x radius=24,y radius=12];
\draw[courseaccent,line width=1pt,->] plot[domain=43:59,samples=12,variable=\a]
 ({24*(cos(\a)-1)},{12*sin(\a)});
\draw[courseaccent,line width=1pt,->] plot[domain=223:239,samples=12,variable=\a]
 ({24*(cos(\a)-1)},{12*sin(\a)});
\draw[aux line,dashed] (-24,-12)--(-24,12);
\fill[courseaccent] (0,0) circle[radius=.7];
\node[anchor=north west,inner sep=1mm] at (0,0) {$0,\ T_0$};
\node[above=1mm,inner sep=0pt] at (-24,12) {$T_0/4$};
\node[below=1mm,inner sep=0pt] at (-24,-12) {$3T_0/4$};
\node[below=1.2mm,inner sep=0pt] at (-48,0) {$T_0/2$};
\draw[aux line] (-48,0)--(-48,14);
\node[above=1mm,inner sep=0pt] at (-46,14) {$-4r_0\Delta\widehat V_r$};
\node[anchor=south west,inner sep=0pt] at (-18,12) {$r_0\Delta\widehat V_r$};
\node[anchor=north west,inner sep=0pt] at (-18,-12) {$-r_0\Delta\widehat V_r$};
\end{tikzpicture}

\end{column}
\begin{column}{.47\textwidth}
{\eqcompact
\setlength{\abovedisplayskip}{3pt}\setlength{\belowdisplayskip}{3pt}
\setlength{\abovedisplayshortskip}{2pt}\setlength{\belowdisplayshortskip}{2pt}
\[
 \left(\frac{\eta+2r_0{\Delta\widehat V_r}}{2r_0{\Delta\widehat V_r}}\right)^2
 +\left(\frac{\Delta r}{r_0{\Delta\widehat V_r}}\right)^2=1.
\]
}
\begin{itemize}
 \setlength{\itemsep}{2pt}
 \item $T_0/4$: наибольшее увеличение радиуса.
 \item $T_0/2$: наибольшее отставание, $\Delta r=0$.
 \item $3T_0/4$: наименьший радиус.
 \item $T_0$: возврат в исходную относительную точку.
\end{itemize}
\end{column}
\end{columns}
\end{frame}

\begin{frame}{Сближение радиальными импульсами}
\relax
\begin{columns}[T,onlytextwidth]
\begin{column}{.50\textwidth}
\centering% Baranov, Sec. 1.4.1, Fig. 1.11(b). Native size: 70 x 40 mm.
% Repeated negative radial impulses: lower half-ellipses A -> B -> C -> O.
% With no impulse at C, passive upper half-ellipse returns to B.
% Equal arcs: Delta V_r < 0 at A and O; 2 Delta V_r at B and C.
% All axes have the same scale; the radial/along-track semiaxis ratio is 1:2.
\begin{tikzpicture}[course diagram,x=1mm,y=1mm]
\path[use as bounding box] (-63,-21) rectangle (7,19);
\draw[vector line] (-61,0)--(6,0) node[above] {$\eta$};
\draw[vector line] (0,-15)--(0,17) node[above right] {$\Delta r$};
\foreach \c in {-47.5,-28.5,-9.5}{
 \draw[orbit line] plot[domain=180:360,samples=61,variable=\a]
  ({\c+9.5*cos(\a)},{4.75*sin(\a)});
 \draw[courseaccent,line width=1pt,->] plot[domain=258:276,samples=9,variable=\a]
  ({\c+9.5*cos(\a)},{4.75*sin(\a)});}
\draw[courseink,line width=.7pt,dashed] plot[domain=0:180,samples=61,variable=\a]
 ({-28.5+9.5*cos(\a)},{4.75*sin(\a)});
\draw[courseink,line width=.7pt,->] plot[domain=80:98,samples=9,variable=\a]
 ({-28.5+9.5*cos(\a)},{4.75*sin(\a)});
\foreach \x/\lab in {-57/A,-38/B,0/O}{
 \fill[courseaccent] (\x,0) circle[radius=.7];
 \node[above left,inner sep=1mm] at (\x,0) {$\lab$};}
\fill[courseaccent] (-19,0) circle[radius=.7];
\node[above right,inner sep=1mm] at (-19,0) {$C$};
\foreach \x/\h/\lab in {-57/4/{\Delta V_r},-38/8/{2\Delta V_r},-19/8/{2\Delta V_r},0/4/{\Delta V_r}}{
 \draw[courseaccent,line width=.9pt,->] (\x,0)--(\x,-\h);
 \ifnum\x=0
  \node[below right=1mm,inner sep=0pt,text=courseaccent] at (\x,-\h) {$\lab$};
 \else
  \node[below=1mm,inner sep=0pt,text=courseaccent] at (\x,-\h) {$\lab$};
 \fi}
\node[text=coursemuted,font=\fontsize{8}{10}\selectfont]
 at (-29,15) {без импульса в $C$: возврат в $B$};
\node[text=coursemuted,font=\fontsize{8}{10}\selectfont]
 at (-28,-17) {Радиальные импульсы, $\Delta V_r<0$};
\end{tikzpicture}

\end{column}
\begin{column}{.47\textwidth}
{\justifying
В точке $A$ КА находится на орбите цели, позади неё. Импульс $\Delta V_r<0$ начинает движение по нижней половине эллипса.\par}
\smallskip
Через полвитка КА возвращается на исходный радиус и приближается к цели на $4r_0|\Delta\widehat V_r|$.
\par\smallskip
{\justifying
В $B,C$ импульсы $2\Delta V_r$ меняют знак радиальной скорости. В $O$ импульс $\Delta V_r$ гасит её.\par}
\end{column}
\end{columns}
\smallskip
{\justifying
\textbf{При пропуске импульса в $C$} КА возвращается в $B$ по верхней половине эллипса.\par}
\end{frame}

\begin{frame}{Трансверсальный импульс: орбита}
\relax
\begin{columns}[T,onlytextwidth]
\begin{column}{.50\textwidth}
\centering% Baranov, Sec. 1.4.2, Fig. 1.12; first-order orbit shape.
% Native size: 70 x 40 mm. Black: reference circle; blue dashed: perturbed orbit.
\begin{tikzpicture}[course diagram,x=1mm,y=1mm]
\path[use as bounding box] (-33,-20) rectangle (37,20);
\draw[vector line] (-25,0)--(26,0) node[above] {$X$};
\draw[vector line] (0,-18)--(0,18) node[left] {$Y$};
\draw[courseink,line width=.65pt] (0,0) circle[radius=12];
\draw[orbit line,dashed] plot[domain=0:360,samples=181,variable=\a]
 ({12*(1+.32*(1-cos(\a)))*cos(\a)},
  {12*(1+.32*(1-cos(\a)))*sin(\a)});
\draw[vector line] (0,0)--(125:12) node[midway,left] {$r_0$};
\draw[courseaccent,line width=1.1pt,->] (12,0)--(12,11)
 node[pos=.75,right=1mm] {$\Delta V_t$};
\draw[courseaccent,line width=1.1pt,->] (0,0)--(7,0)
 node[midway,below=1mm] {$\Delta\boldsymbol e$};
\draw[aux line] (-19.68,0)--(-23,8)--(-28,8);
\node[anchor=south,inner sep=1mm] at (-25,8) {$\Delta r_\alpha$};
\fill[courseaccent] (12,0) circle[radius=.65];
\node[anchor=north west,inner sep=1mm] at (12,0) {$P$};
\node[anchor=west,text=coursemuted,font=\fontsize{8}{10}\selectfont]
 at (-26,-19) {$\Delta r_\alpha=4r_0\Delta\widehat V_t,\quad\Delta r_\pi=0$};
\end{tikzpicture}

\figcaption{Чёрная линия --- опорный круг; голубой пунктир --- орбита после импульса. Отклонения увеличены.}
\end{column}
\begin{column}{.47\textwidth}
Положительный трансверсальный импульс: $\Delta\widehat V_t>0$.
{\eqcompact
\setlength{\abovedisplayskip}{3pt}\setlength{\belowdisplayskip}{3pt}
\setlength{\abovedisplayshortskip}{2pt}\setlength{\belowdisplayshortskip}{2pt}
\[
 \begin{aligned}
 \Delta r&=2r_0{\Delta\widehat V_t}(1-\cos\theta),\\
 \Delta a&=2{\Delta\widehat V_t}.
 \end{aligned}
\]
}
{\justifying
Точка импульса становится перицентром новой орбиты; противоположная точка --- апоцентром.\par}
{\eqcompact
\setlength{\abovedisplayskip}{3pt}\setlength{\belowdisplayskip}{3pt}
\setlength{\abovedisplayshortskip}{2pt}\setlength{\belowdisplayshortskip}{2pt}
\[
 \Delta r_\pi=0,\qquad \Delta r_\alpha=4r_0{\Delta\widehat V_t}.
\]
}
{\justifying
При ${\Delta\widehat V_t}<0$ точка импульса становится апоцентром; радиус противоположной апсиды уменьшается.\par}
\end{column}
\end{columns}
\end{frame}

\begin{frame}{Трансверсальный импульс: фаза}
\relax
\begin{columns}[T,onlytextwidth]
\begin{column}{.50\textwidth}
\centering% Baranov, Sec. 1.4.2, Fig. 1.13. Native size: 70 x 40 mm.
\begin{tikzpicture}[course diagram,x=1mm,y=1mm]
\path[use as bounding box] (-9,-28) rectangle (61,12);
\draw[vector line] (0,-25)--(0,9) node[above right] {$\Delta u$};
\draw[vector line] (0,0)--(57,0) node[above] {$\theta$};
\draw[coursemuted,line width=.55pt,dashed] (0,0)--(50,-21);
\draw[aux line,dashed] (50,0)--(50,-21);
\draw[orbit line] plot[domain=0:360,samples=181,variable=\a]
 ({\a/7.2},{21/(6*pi)*(4*sin(\a)-3*\a*pi/180)});
\foreach \x/\lab in {25/\pi,50/2\pi}{
 \draw[courseink,line width=.5pt] (\x,-.8)--(\x,.8);
 \node[above=1mm,inner sep=0pt] at (\x,0) {$\lab$};}
\node[below=1mm,inner sep=0pt] at (40,-21) {$-6\pi\Delta\widehat V_t$};
\draw[aux line] (34,-9)--(34,-14.28);
\node[above,inner sep=1mm,text=coursemuted] at (34,-9)
 {$-3\theta\Delta\widehat V_t$};
\fill[courseaccent] (50,-21) circle[radius=.65];
\node[below left,inner sep=1mm] at (0,0) {$0$};
\end{tikzpicture}

\end{column}
\begin{column}{.47\textwidth}
Для $\Delta\widehat V_t>0$:
{\eqcompact
\setlength{\abovedisplayskip}{3pt}\setlength{\belowdisplayskip}{3pt}
\setlength{\abovedisplayshortskip}{2pt}\setlength{\belowdisplayshortskip}{2pt}
\[
 \Delta u={\Delta\widehat V_t}(4\sin\theta-3\theta).
\]
}
Максимум опережения:
{\eqcompact
\setlength{\abovedisplayskip}{3pt}\setlength{\belowdisplayskip}{3pt}
\setlength{\abovedisplayshortskip}{2pt}\setlength{\belowdisplayshortskip}{2pt}
\[
 \begin{aligned}
 \theta_{\max}&\approx41^\circ24'35'',\\[3pt]
 \Delta u_{\max}&\approx0{,}4776\Delta\widehat V_t.
 \end{aligned}
\]
}
После $\theta\approx73^\circ05'32''$ начинается отставание.
\smallskip

За один опорный период:
{\eqcompact
\setlength{\abovedisplayskip}{3pt}\setlength{\belowdisplayskip}{3pt}
\setlength{\abovedisplayshortskip}{2pt}\setlength{\belowdisplayshortskip}{2pt}
\[
 \Delta u(2\pi)=-6\pi {\Delta\widehat V_t}.
\]
}
\end{column}
\end{columns}
\end{frame}

\begin{frame}{Трансверсальный импульс: траектория}
\relax
\centering% Baranov, Sec. 1.4.2, Fig. 1.14. Native size: 140 x 36 mm.
% r0 = 6700 km, Delta V_t = 0.002 km/s, mu = 398600.4418 km^3/s^2.
% eta = (Delta V_t/n0)(4 sin theta - 3 theta),
% Delta r = 2(Delta V_t/n0)(1-cos theta); axes both in kilometres.
\begin{tikzpicture}[course diagram,x=1mm,y=1mm]
\path[use as bounding box] (-122,-8) rectangle (18,28);
\pgfmathsetmacro{\impulselength}{1.7372939683809365}
\draw[vector line] (-119,0)--(15,0) node[above] {$\eta$, км};
\draw[vector line] (0,-5)--(0,25) node[above left] {$\Delta r$, км};
\foreach \a in {-40,-30,-20,-10}{
 \draw[courseink,line width=.5pt] ({2.6*\a},-.8)--({2.6*\a},.8);
 \node[below=1.2mm,inner sep=0pt] at ({2.6*\a},0) {$\a$};}
\foreach \a in {3.5,7}{
 \draw[courseink,line width=.5pt] (-.8,{3*\a})--(.8,{3*\a});
 \node[right=1.2mm,inner sep=0pt] at (0,{3*\a}) {$\a$};}
\begin{scope}
\clip (-119,-.3) rectangle (6,24);
\draw[orbit line] plot[domain=0:360,samples=181,variable=\a]
 ({2.6*\impulselength*(4*sin(\a)-3*\a*pi/180)},
  {6*\impulselength*(1-cos(\a))});
\draw[courseaccent,line width=.9pt,dashed] plot[domain=360:550,samples=101,variable=\a]
 ({2.6*\impulselength*(4*sin(\a)-3*\a*pi/180)},
  {6*\impulselength*(1-cos(\a))});
\draw[courseaccent,line width=1pt,->] plot[domain=114:122,samples=9,variable=\a]
 ({2.6*\impulselength*(4*sin(\a)-3*\a*pi/180)},
  {6*\impulselength*(1-cos(\a))});
\end{scope}
\fill[courseaccent] (0,0) circle[radius=.65];
\node[below right,inner sep=1mm] at (0,0) {$0$};
\fill[courseaccent] ({-2.6*\impulselength*6*pi},0) circle[radius=.65];
\node[below=1.2mm,inner sep=0pt] at ({-2.6*\impulselength*6*pi},0) {$T_0$};
\node[text=coursemuted,font=\fontsize{8}{10}\selectfont] at (-103,25)
 {следующий виток};
\end{tikzpicture}
\par
\raggedright
$\Delta V_t=2$ м/с, $r_0=6700$ км.
\smallskip

$\Delta r_{\max}\approx6{,}95$ км;
$\eta(T_0)\approx-32{,}75$ км.
\smallskip

{\justifying
Сначала КА смещается вперёд и наружу, затем отстаёт, продолжая подъём. После апоцентра он снижается до исходного радиуса позади невозмущённого КА. Пунктир --- начало следующего витка.\par}
\end{frame}

\begin{frame}{Нормальный импульс}
\relax
\begin{columns}[T,onlytextwidth]
\begin{column}{.50\textwidth}
\centering% Baranov, Sec. 1.4.3, Fig. 1.16. Native size: 70 x 40 mm.
% Oblique projection; the line of nodes is the x axis.
% The change of radius is second order and is omitted.
\begin{tikzpicture}[course diagram,x=1mm,y=1mm]
\path[use as bounding box] (-32,-20) rectangle (38,20);
\draw[vector line] (0,0)--(-14,-16) node[below left] {$X$};
\draw[vector line] (0,0)--(30,0) node[right] {$Y$};
\draw[vector line] (0,0)--(0,18) node[above left] {$Z$};
\draw[courseink,line width=.65pt] plot[domain=0:360,samples=181,variable=\a]
 ({-10*cos(\a)+23*sin(\a)},{-11*cos(\a)});
\draw[orbit line,dashed] plot[domain=0:360,samples=181,variable=\a]
 ({-10*cos(\a)+23*sin(\a)},{-11*cos(\a)+7*sin(\a)});
\draw[courseaccent,line width=1.1pt,->] (-10,-11)--(-10,-2.5)
 node[pos=.8,left=1mm] {$\Delta V_z$};
\fill[courseaccent] (-10,-11) circle[radius=.7];
\node[below right,inner sep=1mm] at (-10,-11) {$P$};
\draw[aux line] (0,0)--(23,7);
\draw[courseaccent,line width=.7pt] (10,0) arc[start angle=0,end angle=16.93,radius=10];
\node[inner sep=0pt] at (13.1,1.8) {$\Delta i$};
\draw[courseaccent,line width=.65pt,<->] (24.5,0)--(24.5,7)
 node[midway,right=1mm] {$|z|_{\max}$};
\draw[aux line] (23,7)--(26,7);
\node[text=coursemuted,font=\fontsize{8}{10}\selectfont] at (3,-18)
 {$|z|_{\max}=r_0|\Delta\widehat V_z|$};
\end{tikzpicture}

\figcaption{Чёрная линия --- опорный круг; голубой пунктир --- орбита после импульса. Угол поворота увеличен.}
\end{column}
\begin{column}{.47\textwidth}
После нормального импульса $\Delta\widehat V_z$:
{\eqcompact
\setlength{\abovedisplayskip}{3pt}\setlength{\belowdisplayskip}{3pt}
\[
 \begin{aligned}
 z&=r_0{\Delta\widehat V_z}\sin\theta,\\
 V_z&=V_0{\Delta\widehat V_z}\cos\theta.
 \end{aligned}
\]
}
Угол поворота плоскости: $\Delta i=|\Delta\widehat V_z|$.
\smallskip

Амплитуда отклонения: $|z|_{\max}=r_0|{\Delta\widehat V_z}|$.
\smallskip

{\justifying
В первом порядке $\Delta r=\Delta u=0$: движение относительно опорного КА состоит из колебаний по нормали.\par}

\end{column}
\end{columns}
\end{frame}

% Глава 2 методички: основной текст и формулы, без блоков explanation.
\begin{frame}{Постановка задачи оптимизации}
\relax
\derivationmathspacing
\begin{columns}[T,onlytextwidth]
\begin{column}{.49\textwidth}
\begin{block}{Граничные условия}
\justifying
Пусть заданы начальная и конечная орбиты, начальное и конечное
положения аппарата и моменты $t_i,t_f$.
\par\smallskip
Опорная плоскость совпадает с начальной, а ось $X$ направлена
по линии пересечения плоскостей к восходящему узлу конечной орбиты.
\end{block}
\end{column}
\begin{column}{.48\textwidth}
\begin{block}{Параметры импульсов}
\justifying
Сначала считаем число импульсов $N$ выбранным и все импульсы
ненулевыми. Неизвестны их компоненты и фазы $\varphi_j$.
Функционал равен\par
{\centering$\displaystyle J=\sum\nolimits_{j=1}^N\Delta\widehat V_j.$\par}
Крышка над импульсом означает деление на $V_0$.
\end{block}
\end{column}
\end{columns}
Величина $\Delta\widehat V_j$ есть модуль импульса:
{\eqcompact\derivationmathspacing
\[
 \Delta\widehat V_j=
 \sqrt{(\Delta\widehat V_{rj})^2+(\Delta\widehat V_{tj})^2+
                         (\Delta\widehat V_{zj})^2}\geq0.
\]
}
\end{frame}

\begin{frame}{Полная система граничных условий}
\relax
\derivationmathspacing
{\eqcompact\derivationmathspacing
\centering\(\displaystyle
\begin{aligned}[t]
 &2\sum\nolimits_j\Delta\widehat V_{tj}=\Delta a,\\[4pt]
 &\sum\nolimits_j(\Delta\widehat V_{rj}\sin\varphi_j+
              2\Delta\widehat V_{tj}\cos\varphi_j)=\Delta e_x,\\[4pt]
 &\sum\nolimits_j(-\Delta\widehat V_{rj}\cos\varphi_j+
              2\Delta\widehat V_{tj}\sin\varphi_j)=\Delta e_y,\\[4pt]
 &-\sum\nolimits_j\Delta\widehat V_{zj}\sin\varphi_j=0,\\[4pt]
 &\sum\nolimits_j\Delta\widehat V_{zj}\cos\varphi_j=\Delta i,\\[4pt]
 &\sum\nolimits_j\bigl[2\Delta\widehat V_{rj}(1-\cos\varphi_j)
       +\Delta\widehat V_{tj}(-3\varphi_j+4\sin\varphi_j)\bigr]
       =\Delta t-\frac32\varphi_f\Delta a\\[2pt]
 &\hspace{10mm}+2\Delta e_x\sin\varphi_f
                   +2\Delta e_y(1-\cos\varphi_f).
\end{aligned}
\)\par}
\medskip
\justifying
Аргумент широты $\varphi$ отсчитывается от линии узлов.
\end{frame}

\begin{frame}{Функция Лагранжа}
\relax
\derivationmathspacing
\justifying
Введём постоянные множители Лагранжа $l_1,\ldots,l_6$
в порядке строк системы.
{\eqcompact\derivationmathspacing
\setlength{\jot}{1pt}
\[
\begin{aligned}
 \mathcal L={}&J-l_1\left[2\sum\nolimits_j\Delta\widehat V_{tj}-\Delta a\right]\\
 &-l_2\left[\sum\nolimits_j(\Delta\widehat V_{rj}\sin\varphi_j+
                    2\Delta\widehat V_{tj}\cos\varphi_j)-\Delta e_x\right]\\
 &-l_3\left[\sum\nolimits_j(-\Delta\widehat V_{rj}\cos\varphi_j+
                    2\Delta\widehat V_{tj}\sin\varphi_j)-\Delta e_y\right]\\
 &-l_4\left[-\sum\nolimits_j\Delta\widehat V_{zj}\sin\varphi_j\right]\\
 &-l_5\left[\sum\nolimits_j\Delta\widehat V_{zj}\cos\varphi_j
                    -\Delta i\right]\\
 &-l_6\left[\sum\nolimits_j\bigl[2\Delta\widehat V_{rj}(1-\cos\varphi_j)
       +\Delta\widehat V_{tj}(-3\varphi_j+4\sin\varphi_j)\bigr]\right.\\
 &\hspace{20mm}\left.-\Delta t+\frac32\varphi_f\Delta a
         -2\Delta e_x\sin\varphi_f-2\Delta e_y(1-\cos\varphi_f)\right].
\end{aligned}
\]
}
\end{frame}

\begin{frame}{Стационарность по компонентам импульса}
\relax
\derivationmathspacing
Приравняем к нулю производные функции Лагранжа по трём компонентам $j$-го импульса:
{\eqcompact\derivationmathspacing
\[
\begin{aligned}
 \frac{\partial\mathcal L}{\partial\Delta\widehat V_{rj}}
 &=\frac{\Delta\widehat V_{rj}}{\Delta\widehat V_j}
       -l_2\sin\varphi_j+l_3\cos\varphi_j-2l_6(1-\cos\varphi_j)=0,\\[2pt]
 \frac{\partial\mathcal L}{\partial\Delta\widehat V_{tj}}
 &=\frac{\Delta\widehat V_{tj}}{\Delta\widehat V_j}
       -2l_1-2l_2\cos\varphi_j-2l_3\sin\varphi_j
       -l_6(-3\varphi_j+4\sin\varphi_j)=0,\\[2pt]
 \frac{\partial\mathcal L}{\partial\Delta\widehat V_{zj}}
 &=\frac{\Delta\widehat V_{zj}}{\Delta\widehat V_j}
       +l_4\sin\varphi_j-l_5\cos\varphi_j=0.
\end{aligned}
\]
}
\end{frame}

\begin{frame}{Базис-вектор Лоудена}
\relax
\derivationmathspacing
Перейдём к константам Баранова:
{\derivationmathspacing
\[
\boxed{\begin{aligned}
 \lambda_1&=l_1,&\lambda_2&=l_3+2l_6,&\lambda_3&=l_2,\\
 \lambda_4&=-l_4,&\lambda_5&=l_5,&\lambda_6&=l_6.
\end{aligned}}
\]
}{\derivationmathspacing
\[
\boxed{\begin{aligned}
 s_r(\varphi)&=\lambda_3\sin\varphi-\lambda_2\cos\varphi+2\lambda_6,\\
 s_t(\varphi)&=2\lambda_1+2\lambda_3\cos\varphi
                         +2\lambda_2\sin\varphi-3\lambda_6\varphi,\\
 s_z(\varphi)&=\lambda_4\sin\varphi+\lambda_5\cos\varphi.
\end{aligned}}
\]
}
\medskip
\justifying
Вектор $\boldsymbol s=(s_r,s_t,s_z)$ называется базис-вектором Лоудена.
У Баранова те же компоненты обозначаются также $\lambda,\mu,\nu$;
им соответствуют $s_r,s_t,s_z$ в таком же порядке.
\end{frame}

\begin{frame}{Направления и фазы импульсов}
\relax
\derivationmathspacing
{\derivationmathspacing
\[
\Delta\widehat{\boldsymbol V}_j=
           \Delta\widehat V_j\boldsymbol s(\varphi_j),
 \qquad \|\boldsymbol s(\varphi_j)\|=1.
\]
}
\justifying
Для импульса, расположенного внутри рассматриваемого интервала,
можно варьировать также фазу $\varphi_j$.
Поэтому необходимое условие стационарности имеет вид
{\derivationmathspacing
\[
\frac{\partial\mathcal L}{\partial\varphi_j}
 =-\boldsymbol s'(\varphi_j)\cdot\Delta\widehat{\boldsymbol V}_j=0,
 \qquad
 \left.\frac{\mathrm d}{\mathrm d\varphi}\|\boldsymbol s(\varphi)\|^2
                  \right|_{\varphi_j}=0.
\]
}
\medskip
\justifying
Обычное правило множителей Лагранжа даёт необходимые условия
при независимости дифференциалов ограничений по свободным
параметрам манёвра и ненулевых импульсах.
Одна стационарность ещё не доказывает минимум и не исключает
улучшения схемы изменением числа импульсов.
\end{frame}

\begin{frame}{Базис-вектор со сдвигом фазы}
\relax
\derivationmathspacing
При $\lambda_2^2+\lambda_3^2>0$ определим угол
{\eqcompact\derivationmathspacing
\setlength{\jot}{0pt}
\[
\cos\alpha=\frac{\lambda_3}{\sqrt{\lambda_2^2+\lambda_3^2}},
 \qquad
 \sin\alpha=\frac{\lambda_2}{\sqrt{\lambda_2^2+\lambda_3^2}}.
\]
}
Тогда полный базис-вектор принимает вид
{\eqcompact\derivationmathspacing
\setlength{\jot}{0pt}
\[
\boxed{\begin{aligned}
 s_r&=\sqrt{\lambda_2^2+\lambda_3^2}\,\sin(\varphi-\alpha)+2\lambda_6,\\
 s_t&=2\lambda_1+2\sqrt{\lambda_2^2+\lambda_3^2}\,
                       \cos(\varphi-\alpha)-3\lambda_6\varphi,\\
 s_z&=\frac{\lambda_3\lambda_4-\lambda_2\lambda_5}
                {\sqrt{\lambda_2^2+\lambda_3^2}}\,\sin(\varphi-\alpha)+\frac{\lambda_2\lambda_4+\lambda_3\lambda_5}
                {\sqrt{\lambda_2^2+\lambda_3^2}}\,\cos(\varphi-\alpha).
\end{aligned}}
\]
}
\end{frame}

\begin{frame}{Взвешенная сумма граничных условий}
\relax
\derivationmathspacing
\justifying
Умножим шесть граничных уравнений на $l_1,\ldots,l_6$ и сложим.
Коэффициенты при компонентах каждого импульса совпадают
с правыми частями условий на направляющие косинусы, поэтому
{\derivationmathspacing
\[
\begin{aligned}
 \sum\nolimits_j\boldsymbol s(\varphi_j)\cdot\Delta\widehat{\boldsymbol V}_j
   ={}&l_1\Delta a+l_2\Delta e_x+l_3\Delta e_y+l_5\Delta i\\
      &+l_6\left[\Delta t-\frac32\varphi_f\Delta a
          +2\Delta e_x\sin\varphi_f+2\Delta e_y(1-\cos\varphi_f)\right].
\end{aligned}
\]
}
\medskip
\justifying
Правая часть зависит только от заданных конечных условий
и выбранных множителей. Для любого набора импульсов,
выполняющего эти условия, она одинакова.
\end{frame}

\begin{frame}{Добавление нового импульса}
\relax
\derivationmathspacing
\justifying
Теперь допускаем изменение числа импульсов. В произвольной
точке $\varphi$ добавим малый импульс
\[
 \Delta\widehat{\boldsymbol V}_{\text{нов}}=\eta\boldsymbol d,
 \qquad \eta>0,\qquad \|\boldsymbol d\|=1.
\]
Здесь $\eta$ --- его безразмерный модуль, а $\boldsymbol d$ ---
единичный вектор направления. Одновременно корректируем прежние
импульсы и их свободные фазы, сохраняя все шесть граничных условий.
Возможность такой малой компенсации обеспечивает уже принятое
условие регулярности ограничений.
\par\smallskip
Множители $l_1,\ldots,l_6$ при сравнении сохраняем теми же,
что у исходного манёвра. У оптимального манёвра такая допустимая
поправка не должна уменьшать $J$.
{\derivationmathspacing
\[
\|\boldsymbol s(\varphi)\|\leq1
 \qquad\text{при всех допустимых }\varphi\in[\varphi_i,\varphi_f].
\]
}
\end{frame}

\begin{frame}{Достаточные условия оптимальности}
\relax
\derivationmathspacing
\justifying
Теперь докажем достаточность. Для любого набора множителей,
обеспечивающего ограничение $\|\boldsymbol s(\varphi)\|\leq1$,
получается нижняя оценка затрат:
{\derivationmathspacing
\[
\begin{aligned}
 J\geq{}&l_1\Delta a+l_2\Delta e_x+l_3\Delta e_y+l_5\Delta i\\
       &+l_6\left[\Delta t-\frac32\varphi_f\Delta a
          +2\Delta e_x\sin\varphi_f+2\Delta e_y(1-\cos\varphi_f)\right].
\end{aligned}
\]
}
\medskip
\justifying
Таким образом, при
$\Delta\widehat{\boldsymbol V}_j=\Delta\widehat V_j\boldsymbol s(\varphi_j)$
нижняя оценка достигается, и никакой другой допустимый манёвр
не имеет меньших затрат.
\par\smallskip
При $\lambda_6\ne0$ базис-вектор непериодичен из-за члена
$-3\lambda_6\varphi$, поэтому проверять его только на одном
произвольно выбранном обороте недостаточно.
\end{frame}

\begin{frame}{Достаточные условия глобального минимума}
\relax
\derivationmathspacing
\justifying
В линейной задаче с произвольным числом импульсов манёвр
обеспечивает глобальный минимум $J$, если выполнены следующие условия.
\begin{enumerate}
\setlength{\itemsep}{0pt}
 \item Импульсы выполняют все шесть граничных
 условий, а их фазы принадлежат
 допустимому интервалу $[\varphi_i,\varphi_f]$.
 \item Существует один набор постоянных
 $\lambda_1,\ldots,\lambda_6$, задающий базис-вектор
 по формулам слайда~22, для которого в точке
 каждого ненулевого импульса выполнено:
 {\eqcompact\derivationmathspacing
\[
\boldsymbol s(\varphi_j)
     =\frac{\Delta\widehat{\boldsymbol V}_j}{\Delta\widehat V_j},
  \qquad \|\boldsymbol s(\varphi_j)\|=1.
\]
}
 \item Для того же базис-вектора выполнено достаточное
 условие:
 {\eqcompact\derivationmathspacing
\[
\|\boldsymbol s(\varphi)\|\leq1
  \qquad\text{при всех }\varphi\in[\varphi_i,\varphi_f].
\]
}
\end{enumerate}
\end{frame}

\begin{frame}{Достаточные условия глобального минимума}
\relax
\derivationmathspacing
\justifying
Внутри интервала в точках ненулевых импульсов также выполнено
условие стационарности:
{\eqcompact\derivationmathspacing
\[
\left.\frac{\mathrm d}{\mathrm d\varphi}\|\boldsymbol s(\varphi)\|^2
          \right|_{\varphi_j}=0.
\]
}
Оно следует из второго и третьего пунктов: в точке импульса
квадрат модуля достигает своего максимального значения $1$.
На фиксированных концах интервала равенство производной нулю
не требуется.

Перечисленная совокупность \emph{достаточна} для глобального
минимума без предположения регулярности. При принятой выше
регулярности она также \emph{необходима}: для оптимального
манёвра существуют множители, удовлетворяющие этим условиям.
Таким образом, для регулярного случая получены необходимые
и достаточные условия оптимальности.
\end{frame}

\begin{frame}{Переход на заданную орбиту}
\relax
\derivationmathspacing
\justifying
Конечная орбита задана, а положение аппарата на ней и длительность
перехода свободны. Фазовая строка больше не ограничивает
импульсы. Соответствующий множитель исчезает:
$l_6=0$, $\lambda_6=0$.
{\derivationmathspacing
\[
\boxed{\begin{aligned}
 s_r(\varphi)&=\lambda_3\sin\varphi-\lambda_2\cos\varphi,\\
 s_t(\varphi)&=2\lambda_1+2\lambda_3\cos\varphi+2\lambda_2\sin\varphi,\\
 s_z(\varphi)&=\lambda_4\sin\varphi+\lambda_5\cos\varphi.
\end{aligned}}
\]
}
\medskip
\justifying
Условие на плоскость орбиты не снято, поэтому
$\lambda_4,\lambda_5$ в общем случае ненулевые.
\par\smallskip
При задаче встречи фазовое условие сохраняется;
свободная длительность сама по себе не позволяет исключить $\lambda_6$.
\end{frame}

\begin{frame}{Компланарный переход}
\relax
\derivationmathspacing
\justifying
Если обе орбиты лежат в одной плоскости, $\Delta i=0$.
В линейной модели нормальные импульсы не меняют $a,e_x,e_y$,
а их добавление к заданным плоским компонентам увеличивает
модуль импульса. Следовательно, оптимальное решение можно
искать при $\Delta\widehat V_{zj}=0$.
{\derivationmathspacing
\[
\lambda_6=0,\qquad \lambda_4=\lambda_5=0,\qquad s_z=0.
\]
}
Для $\lambda_2^2+\lambda_3^2>0$ базис-вектор принимает вид
{\derivationmathspacing
\[
\boxed{\begin{aligned}
 s_r&=\sqrt{\lambda_2^2+\lambda_3^2}\,\sin(\varphi-\alpha),\\
 s_t&=2\lambda_1+2\sqrt{\lambda_2^2+\lambda_3^2}\,
                         \cos(\varphi-\alpha),\\
 s_z&=0.
\end{aligned}}
\]
}
\end{frame}

\begin{frame}{Источники}
\relax
\setbeamertemplate{bibliography item}{\insertbiblabel}
\begin{thebibliography}{9}
\bibitem{Baranov}
Баранов А. А. Маневрирование космических аппаратов в окрестности круговой орбиты.
М.: Спутник+, 2016. 511~с. Гл.~1--2.
\bibitem{EdelbaumGobetzWashington1966}
Edelbaum T. N., Gobetz F. W., Washington M.
\href{https://ntrs.nasa.gov/citations/19660030604}{Minimum-Impulse Time-Free Transfer Between Elliptic Orbits}.
NASA CR-636, 1966.
\bibitem{Edelbaum1966}
Edelbaum T. N.
\href{https://doi.org/10.2514/6.1966-7}{How many impulses?}
AIAA Paper 66-7, 1966.
\bibitem{Neustadt1964}
Neustadt L. W.
\href{https://doi.org/10.1137/0302004}{Optimization, a Moment Problem, and Nonlinear Programming}.
Journal of the Society for Industrial and Applied Mathematics, Series A: Control.
1964. Vol.~2, No.~1. P.~33--53.
\end{thebibliography}
\end{frame}

\begin{frame}{Источники}
\relax
\setbeamertemplate{bibliography item}{\insertbiblabel}
\begin{thebibliography}{9}
\setcounter{enumiv}{4}
\bibitem{SternPotter1965}
Stern R. G., Potter J. E.
\href{https://www.sciencedirect.com/science/article/abs/pii/S1474667017690700}{Optimization of Midcourse Velocity Corrections}.
IFAC Proceedings Volumes. 1965. Vol.~2, No.~1. P.~70--83.
\bibitem{Lawden}
Lawden D. F. Optimal Trajectories for Space Navigation.
London: Butterworths, 1963. Гл.~3, \S~5.1, гл.~6.
\bibitem{Edelbaum1967}
Edelbaum T. N.
\href{https://ntrs.nasa.gov/citations/19670020044}{A General Solution for Minimum Impulse Transfers in the Near Vicinity of a Circular Orbit}.
First Compilation of Papers on Trajectory Analysis and Guidance Theory.
NASA SP-141, 1967. P.~37--71.
\bibitem{BaranovDissertation}
Баранов А. А. Разработка методов расчета параметров маневров космических аппаратов
в окрестностях круговой орбиты: дис.\ д-ра физ.-мат.\ наук. М., 2019. 304~с.
\S~1.3, с.~38--39.
\end{thebibliography}
\end{frame}

\end{document}
